\newcommand{\cfg}[1]{\texttt{\seqsplit{#1}}}

The primary flight software running on the Baseboard's STM32H723VGT6 microcontroller is architected as a set of concurrent FreeRTOS tasks. The central task implements a deterministic finite state machine that governs all flight phases, from pre-flight sensor calibration through powered ascent, recovery deployment, and landing. The state machine executes at a fixed rate of \textbf{100\,Hz} (10\,ms period). Each state defines an \emph{entry action}, executed once upon transition, and a \emph{handler} that evaluates transition conditions on every cycle and returns the next state.

\paragraph{State Machine Overview}

The flight state machine comprises 11~states organized into four functional groups: ground operations, powered and ballistic flight, recovery, and terminal states. Figure~\ref{fig:state_machine} illustrates the complete state diagram with all transition conditions.

\begin{figure}[H]
\centering
\resizebox{!}{0.85\textheight}{%
\begin{tikzpicture}[
    state/.style={draw, rounded corners=6pt, thick, minimum width=2.6cm, minimum height=1.2cm, align=center, font=\sffamily\small\bfseries},
    abortst/.style={state, draw=red!70!black, fill=red!8},
    pyrost/.style={state, draw=orange!80!black, fill=orange!8},
    groundst/.style={state, draw=EtseNavy!80, fill=EtseNavy!6},
    flightst/.style={state, draw=green!50!black, fill=green!6},
    terminalst/.style={state, draw=gray!60, fill=gray!8},
    arr/.style={-{Stealth[length=2.5mm]}, thick},
    faultarr/.style={arr, red!60!black},
    cmd/.style={arr, red!60!black, dashed},
    lbl/.style={font=\sffamily\scriptsize, align=center, fill=white, inner sep=1.5pt},
]

% --- Nominal flow: vertical column ---
\node[groundst]   at (0, 0)     (idle)   {1. Idle};
\node[groundst]   at (0, -2.4)  (cal)    {2. Calibration};
\node[groundst]   at (0, -4.8)  (pre)    {3. Pre-launch};
\node[flightst]   at (0, -7.2)  (boost)  {4. Boost};
\node[flightst]   at (0, -9.6)  (coast)  {5. Coast};
\node[flightst]   at (0, -12.0) (actrl)  {6. Active\\Control};
\node[pyrost]     at (0, -14.4) (apogee) {7. Apogee\\{\scriptsize\textit{Drogue}}};
\node[pyrost]     at (0, -16.8) (main)   {8. Main\\Parachute};
\node[terminalst] at (0, -19.2) (landed) {9. Landed};

% --- Abort states (right column) ---
\node[abortst] at (6, -4.8)  (gabort) {10. Ground\\Abort};
\node[abortst] at (6, -14.4) (dabort) {11. Descent\\Abort};

% --- Nominal transitions (black solid, vertical) ---
\draw[arr] (idle)   -- node[lbl, right] {Auto-start} (cal);
\draw[arr] (cal)    -- node[lbl, right] {Sensors\\calibrated} (pre);
\draw[arr] (pre)    -- node[lbl, right] {$a_Y >$ thr.\\$\times N$ samples} (boost);
\draw[arr] (boost)  -- node[lbl, right] {$a_Y <$ thr.\\$\times N$ samples} (coast);
\draw[arr] (coast)  -- node[lbl, right] {Alt $\geq$ thr.} (actrl);
\draw[arr] (actrl)  -- node[lbl, right] {Apogee\\detected} (apogee);
\draw[arr] (apogee) -- node[lbl, right] {Alt $\leq$ thr.\\or timeout} (main);
\draw[arr] (main)   -- node[lbl, right] {Landing\\detected} (landed);

% --- Ground abort (red solid, from ground states to abort) ---
\draw[faultarr] (idle.east)  -- node[lbl, above] {Fault} (gabort.west);
\draw[faultarr] (cal.east)   -- node[lbl, above] {Fault} (gabort.west);
\draw[faultarr] (pre.east)   -- node[lbl, above] {Fault} (gabort.west);

% --- Ground abort reset (dashed, routed back up) ---
\draw[cmd] (gabort.north) |- node[lbl, above, pos=0.8] {CMD Reset} (idle.east);

% --- Descent abort (red solid, from Apogee) ---
\draw[faultarr] (apogee.east) -- node[lbl, above] {Serious sensor\\fault} (dabort.west);

% --- Descent abort to landed (black solid) ---
\draw[arr] (dabort.south) |- node[lbl, above, pos=0.2] {Landing\\detected} (landed.east);

\end{tikzpicture}%
}
\caption{Flight controller state diagram. Black solid arrows: automatic transitions. Red solid arrows: fault-triggered transitions. Red dashed arrows: remote operator command overrides.}
\label{fig:state_machine}
\end{figure}

\paragraph{Ground States}

\textbf{Idle} is the initial state after power-on or reset. On entry, all pyrotechnic channels are disarmed, the calibration context is cleared, and the SD card is mounted and opened for logging. Sensors operate in low-power mode. The transition to \emph{Calibration} occurs automatically after sensor initialization completes when the auto-start flag is enabled (\cfg{AUTO\_START\_CALIBRATION}).

In the \textbf{Calibration} state, SD logging is enabled and all sensors switch to high-performance mode. Three calibration processes execute concurrently:
\begin{itemize}
    \item \textbf{Reference pressure:} The first \cfg{PRESSURE\_CALIBRATION\_DISCARD\_SAMPLES} samples are discarded to allow sensor stabilization; the subsequent \cfg{PRESSURE\_CALIBRATION\_SAMPLES} samples are averaged to establish the ground-level barometric reference.
    \item \textbf{Gyroscope bias:} The first \cfg{GYRO\_CALIBRATION\_DISCARD\_SAMPLES} samples are discarded; the subsequent \cfg{GYRO\_CALIBRATION\_SAMPLES} samples are averaged to determine the zero-rate bias on all three axes.
    \item \textbf{GPS fix:} A valid fix with at least \cfg{GPS\_FIX\_MIN\_SATELLITES} satellite, non-zero position, and non-zero altitude is required.
\end{itemize}
Additionally, the state waits for the onboard Kalman filter to converge before considering calibration complete. The state machine advances to \emph{Pre-launch} only when all three calibrations are simultaneously valid and the filter has converged.

In the \textbf{Pre-launch} state, the rocket is on the launch rail, fully calibrated and awaiting ignition. A confirmation counter is initialized. The transition to \emph{Boost} is triggered when the Y-axis acceleration exceeds the configured threshold (\cfg{PRELAUNCH\_BOOST\_ACCEL\_Y\_THRESHOLD}) for a configurable number of consecutive samples (\cfg{PRELAUNCH\_BOOST\_CONSECUTIVE\_SAMPLES}), indicating motor ignition.

\paragraph{Flight States}

During the \textbf{Boost} state the motor is active and the vehicle experiences sustained positive acceleration. The transition to \emph{Coast} occurs when the measured acceleration drops below the configured threshold (\cfg{BOOST\_COAST\_ACCEL\_Y\_THRESHOLD}) for the required number of consecutive samples (\cfg{BOOST\_COAST\_CONSECUTIVE\_SAMPLES}), indicating motor burnout.

The \textbf{Coast} state represents the unpowered ascent phase after motor burnout. The transition to \emph{Active Control} occurs when the barometric altitude exceeds a configurable height threshold (\cfg{COAST\_ACTIVE\_CONTROL\_BAROM\_ALT\_THRESHOLD}).

During the \textbf{Active Control} phase the vehicle is in ascending ballistic flight. The active airbrake control algorithm (detailed in Section~\ref{subsec:active_flight_control}) modulates drag to converge on the target apogee. Multiple apogee detection criteria are monitored redundantly in a logical OR configuration, any single criterion triggers the transition to \emph{Apogee} (EuRoC-LV-RQT-0350). Each criterion can be individually enabled or disabled via its corresponding configuration flag:
\begin{itemize}
    \item Barometric altitude $\geq$ threshold (\cfg{ACTIVE\_CONTROL\_APOGEE\_BAROM\_ALT\_THRESHOLD}).
    \item GPS altitude $\geq$ threshold (\cfg{ACTIVE\_CONTROL\_APOGEE\_GPS\_ALT\_THRESHOLD}).
    \item GPS vertical velocity $\leq$ threshold (\cfg{ACTIVE\_CONTROL\_APOGEE\_GPS\_VEL\_Y\_THRESHOLD}), indicating the onset of descent.
    \item Safety timeout from state entry (\cfg{ACTIVE\_CONTROL\_APOGEE\_DELAY\_MS}, enabled via \cfg{ACTIVE\_CONTROL\_APOGEE\_DELAY\_ENABLED}).
\end{itemize}

\paragraph{Recovery States}

On entry to the \textbf{Apogee} state, the drogue parachute pyrotechnic channel is immediately fired via the Powerboard relay driver (see Section~\ref{subsubsection:avionics_bay}). The transition to \emph{Main Parachute} occurs when either the barometric altitude descends to the configured threshold (\cfg{APOGEE\_MAIN\_PARACHUTE\_BAROM\_ALT\_THRESHOLD}) or the configured timeout elapses (\cfg{APOGEE\_MAIN\_PARACHUTE\_DELAY\_MS}, enabled via \cfg{APOGEE\_MAIN\_PARACHUTE\_DELAY\_ENABLED}).

On entry to the \textbf{Main Parachute} state, the main parachute pyrotechnic channel is fired. Per EuRoC competition rules, the main parachute must not be deployed at an altitude higher than 450\,m AGL (EuRoC-LV-RQT-0370); this limit is enforced through the configurable threshold (\cfg{APOGEE\_MAIN\_PARACHUTE\_BAROM\_ALT\_THRESHOLD}). A landing confirmation counter is initialized. The transition to \emph{Landed} requires both conditions to be satisfied simultaneously for the configured number of consecutive samples (\cfg{MAIN\_PARACHUTE\_LANDED\_CONSECUTIVE\_SAMPLES}):
\begin{itemize}
    \item Barometric altitude $\leq$ threshold (\cfg{MAIN\_PARACHUTE\_LANDED\_BAROM\_ALT\_THRESHOLD}).
    \item Barometric vertical velocity $\leq$ threshold (\cfg{MAIN\_PARACHUTE\_LANDED\_BAROM\_VEL\_Y\_THRESHOLD}).
\end{itemize}

\paragraph{Terminal and Abort States}

On entry to the \textbf{Landed} state, all pyrotechnic channels are disarmed. After the configured delay (\cfg{LANDED\_SD\_STOP\_DELAY\_MS}), SD logging is disabled and the log file is closed to ensure data integrity. The system remains in this state, continuing to transmit telemetry with GPS position data to facilitate localization of the vehicle.

The \textbf{Ground Abort} state is activated by system fault detection (sensor initialization failures or SD card errors) during any ground state. On entry, all pyrotechnic channels are disarmed, SD logging is stopped, and the file is closed. The operator can recover the system from this state by issuing a remote reset command, which forces a return to \emph{Idle}.

The \textbf{Descent Abort} state is entered when a serious sensor fault is detected during descent. Unlike the nominal recovery sequence, this state skips the main parachute deployment, allowing the vehicle to descend under the drogue parachute alone at a higher descent rate. The flight controller can continue to operate on descent with only the barometer if needed. The state monitors barometric altitude and velocity and transitions to \emph{Landed} upon landing detection.

\paragraph{Redundant Flight Computer}

A secondary flight computer, identical in architecture to the primary Baseboard, is connected to the main controller via a 2-pin signalling interface. The redundant computer carries its own independent set of sensors (barometer, IMU, and magnetometer) and runs the same flight state machine software, adapted for the redundant hardware. Both computers have physical access to the active control servos; the signalling interface determines which computer has authority.

The 2-pin interface encodes a 2-bit status word driven by the primary controller. Table~\ref{tab:redundancy_signals} defines the signalling protocol.

\begin{table}[H]
\centering
\caption{Redundant flight computer signalling protocol.}
\label{tab:redundancy_signals}
\begin{tabularx}{\textwidth}{c l X}
\toprule
\textbf{Signal} & \textbf{State} & \textbf{Effect} \\
\midrule
\texttt{11} & Nominal & Primary has servo authority. Redundant computer tracks state independently but does not actuate. \\
\addlinespace
\texttt{00} & Fault / Loss of connection & Handover: redundant computer assumes servo authority and continues the flight autonomously. \\
\bottomrule
\end{tabularx}
\end{table}

Under nominal operation, the primary controller holds both pins high (\texttt{11}). If the primary detects a serious sensor fault or becomes unresponsive (connection loss), the signal defaults to \texttt{00}, triggering automatic handover. The redundant computer, running the state machine independently using its own sensor suite, then assumes active control of the airbrake servos for the remainder of the flight.

Separately, a COTS CATS Vega flight computer (Section~\ref{subsubsection:avionics_bay}) runs its own open-source firmware with independent apogee detection and has dedicated access to the recovery pyrotechnic channels. Per EuRoC rules, the CATS Vega firmware must remain unmodified. It operates fully decoupled from the SRAD flight software and serves as a last-resort recovery backup.

\paragraph{Confirmation Counter Mechanism}

All safety-critical transitions that rely on sensor thresholds employ a consecutive sample confirmation mechanism to prevent false triggers from transient noise or anomalous readings. A condition must be continuously satisfied for $N$ consecutive state machine cycles before the transition fires. If the condition is violated at any point during the count, the counter resets to zero. Table~\ref{tab:confirm_counters} lists all guarded transitions with their corresponding configuration constants.

\begin{table}[H]
\centering
\caption{Confirmation counter parameters for safety-critical transitions.}
\label{tab:confirm_counters}
\small
\begin{tabularx}{\textwidth}{l l X}
\toprule
\textbf{Transition} & \textbf{Samples} & \textbf{Condition} \\
\midrule
Pre-launch $\rightarrow$ Boost
    & \cfg{..\_CONSECUTIVE\_SAMPLES}
    & $a_Y >$ \cfg{PRELAUNCH\_BOOST\_ACCEL\_Y\_THRESHOLD} \\
\addlinespace
Boost $\rightarrow$ Coast
    & \cfg{..\_CONSECUTIVE\_SAMPLES}
    & $a_Y <$ \cfg{BOOST\_COAST\_ACCEL\_Y\_THRESHOLD} \\
\addlinespace
Main Par. $\rightarrow$ Landed
    & \cfg{..\_CONSECUTIVE\_SAMPLES}
    & Alt $\leq$ \cfg{..\_BAROM\_ALT\_THRESHOLD} $\wedge$ $|v| \leq$ \cfg{..\_VEL\_Y\_THRESHOLD} \\
\addlinespace
Descent Ab. $\rightarrow$ Landed
    & \cfg{..\_CONSECUTIVE\_SAMPLES}
    & Alt $\leq$ \cfg{..\_BAROM\_ALT\_THRESHOLD} $\wedge$ $|v| \leq$ \cfg{..\_VEL\_Y\_THRESHOLD} \\
\bottomrule
\end{tabularx}
\end{table}

\paragraph{Remote Command Interface}

The ground station operator can override automatic transitions at any point during the flight via commands transmitted over the LoRa telemetry link. Commands are processed with higher priority than the state handler: if a valid command is received in a given cycle, the corresponding transition is executed without evaluating the current state's handler logic.

\begin{table}[H]
\centering
\caption{Remote operator commands and their effects on the state machine.}
\label{tab:remote_commands}
\begin{tabularx}{\textwidth}{l X}
\toprule
\textbf{Command} & \textbf{Effect} \\
\midrule
Reset & Force return to \emph{Idle} from any state. \\
\bottomrule
\end{tabularx}
\end{table}

\paragraph{Sensor Power Management}

To optimize power consumption, the flight software dynamically reconfigures the sensor operating modes based on the current state. Sensors are placed in low-power idle mode during ground states prior to calibration and after landing or abort, and switched to high-performance mode for the active flight envelope.

\begin{table}[H]
\centering
\caption{Sensor operating mode allocation by flight state.}
\label{tab:sensor_modes}
\begin{tabularx}{\textwidth}{l l X}
\toprule
\textbf{State Range} & \textbf{Sensor Mode} & \textbf{Sampling Timers} \\
\midrule
Idle & Low power (idle) & Stopped \\
Calibration through Main Parachute & High performance & Active (100\,Hz) \\
Landed / Ground Abort / Descent Abort & Low power (idle) & Stopped \\
\bottomrule
\end{tabularx}
\end{table}

\paragraph{Data Logging}

SD card logging is enabled upon entry to the \emph{Calibration} state and remains active through all subsequent flight states until landing or abort. Each state machine cycle produces a 104-byte record containing raw sensor readings, computed state estimates, GPS data, system fault flags, and the current state identifier. Records are enqueued into a FreeRTOS queue and written asynchronously to the SD card by a dedicated logging task in buffers of \cfg{SD\_LOGGING\_RECORDS\_PER\_BUFFER} records, ensuring that file I/O latency does not affect the determinism of the main state machine loop.
